% !PRINTED 507
\weberpage{507}

Um auch $\varkappa(\omega)$ selbst als Klasseninvariante auffassen zu
können, wendet man die \gesperrt{Gausssche} Transformation (\S\,9, \S\,32)
an. Danach ist:

\[ \varkappa^2\left(\dfrac{\omega}{2}\right) = \dfrac{4\varkappa(\omega)}{[1+\varkappa(\omega)]^2}, \]

woraus:

\[ \varkappa(\omega) = \dfrac{\varkappa^2\left(\dfrac{\omega}{2}\right)[1+\varkappa^2(\omega)]}{4-2\varkappa^2\left(\dfrac{\omega}{2}\right)}. \]

Es ist also $\varkappa(\omega)$ rational ausgedrückt durch $\varkappa^2(\omega)$ und
$\varkappa^2\left(\dfrac{\omega}{2}\right)$, und $\omega' = \dfrac{\omega}{2}$ genügt der Gleichung

\weq{8}{4A\omega'^2 + 2B\omega' + C = 0.}

Ist $C$ eine gerade Zahl, was wir annehmen können, wenn
$D\equiv 0\ (\mathrm{mod}\ 4)$ oder $\equiv 1\ (\mathrm{mod}\ 8)$ ist, so kann in (8) der
Faktor 2 weggehoben werden und die Diskriminante von (8) ist
gleich $D$. Ist aber $C$ ungerade, was bei $D\equiv 5\ (\mathrm{mod}\ 8)$ notwendig ist, so ist die Diskriminante von (8) gleich $4D$. Daraus
folgt nach 2. und 3.:

4. a) Ist $D\equiv 0\ (\mathrm{mod}\ 4)$, so ist $\varkappa(\omega)$ \gesperrt{rational ausdrückbar durch die Klasseninvarianten der Diskriminante $4D$.}

b) Ist $D\equiv 1\ (\mathrm{mod}\ 8)$, so ist $\varkappa(\omega)$ \gesperrt{ausdrückbar durch $\sqrt{D}$ und durch die Klasseninvarianten der Diskriminante $4D$.}

c) Ist $D\equiv 5\ (\mathrm{mod}\ 8)$, so ist $\varkappa(\omega)$ \gesperrt{rational ausdrückbar durch $\sqrt{D}$ und durch die Klasseninvarianten der Diskriminante $16D$.}

\begin{center}
\S\,136.\quad \gesperrt{Quadratische Transformationsgrade.}
\end{center}

Wenn der Transformationsgrad $n$ eine \gesperrt{ungerade Quadratzahl}
ist, so ist nach \S\,72, 5. für ein variables $\omega$:

\weq{1}{M = (-1)^{\frac{a-1}{2}} \partial^3\left(\dfrac{\eta\left(\dfrac{c+\partial\omega}{a}\right)}{\eta(\omega)}\right)^6,}
%CONT>>
